% GENERATED FILE. Edit canonical lessons, then run npm run generate:books.
% canonical-source-hash: fnv1a64:f7efc3eaad26f560
% canonical-lessons: KA-C278-sala-kodu, KA-C278-hintirugisu, KA-C278-hintirugu, KA-C278-tade, KA-C278-tayarisu

\chapter{To lend, To give back, To come back, To prevent, To prepare}
\label{ch:ka-to-lend-to-give-back-to-come-back-to-prevent-to-prepare}
\hlchaptermodality{278}

\begin{chapteropening}
\textbf{By the end of this chapter:} \emph{I can say to lend, to give back, to come back, to prevent and to prepare.}

Complete the last lesson of chapter 278: I can say to lend, to give back, to come back, to prevent and to prepare.
\end{chapteropening}

\section[\texorpdfstring{\kn{ಸಾಲ} \kn{ಕೊಡು} (sāla koḍu)}{ಸಾಲ ಕೊಡು (sāla koḍu)}]{\kn{ಸಾಲ} \kn{ಕೊಡು} (sāla koḍu) — to lend (money)}

\label{lesson:KA-C278-sala-kodu}



\begin{quote}
Before the new one: say the Kannada for to paint, then the Kannada for to download.
\end{quote}

\subsection*{You'll want to know: \kn{ಸಾಲ} \kn{ಕೊಡು}}
\textbf{\kn{ಸಾಲ} \kn{ಕೊಡು}} — \emph{sāla koḍu} — \textquotedblleft{}to lend (money)\textquotedblright{}.

\begin{grammarlens}[title={What you've built}]
One more word for this chapter.
\end{grammarlens}

\subsection*{Your turn}
\begin{itemize}
\raggedright
  \item \emph{Say it:} \emph{sāla koḍu}
  \item \emph{Say it:} \emph{sāla koḍu}, once more
  \item \emph{Say it:} say \emph{sāla koḍu}
  \item \emph{Recall:} say the Kannada for to paint, then the Kannada for to download, then say \emph{sāla koḍu} again
\end{itemize}

\begin{tcolorbox}[breakable,colback=teal!4,colframe=teal!35!black,title={Before you move on}]
What does \kn{ಸಾಲ} \kn{ಕೊಡು} mean? (\textquotedblleft{}To lend (money)\textquotedblright{}.) Say it once more.
\end{tcolorbox}
\section[\texorpdfstring{\kn{ಹಿಂತಿರುಗಿಸು} (hintirugisu)}{ಹಿಂತಿರುಗಿಸು (hintirugisu)}]{\kn{ಹಿಂತಿರುಗಿಸು} (hintirugisu) — to give back, to return (something)}

\label{lesson:KA-C278-hintirugisu}



\begin{quote}
Before the new one: say the Kannada for to download, then the Kannada for to lend.
\end{quote}

\subsection*{You'll want to know: \kn{ಹಿಂತಿರುಗಿಸು}}
\textbf{\kn{ಹಿಂತಿರುಗಿಸು}} — \emph{hintirugisu} — \textquotedblleft{}to give back, to return (something)\textquotedblright{}.

\begin{grammarlens}[title={What you've built}]
One more word for this chapter.
\end{grammarlens}

\subsection*{Your turn}
\begin{itemize}
\raggedright
  \item \emph{Say it:} \emph{hintirugisu}
  \item \emph{Say it:} \emph{hintirugisu}, once more
  \item \emph{Say it:} say \emph{hintirugisu}
  \item \emph{Recall:} say the Kannada for to download, then the Kannada for to lend, then say \emph{hintirugisu} again
\end{itemize}

\begin{tcolorbox}[breakable,colback=teal!4,colframe=teal!35!black,title={Before you move on}]
What does \kn{ಹಿಂತಿರುಗಿಸು} mean? (\textquotedblleft{}To give back, to return (something)\textquotedblright{}.) Say it once more.
\end{tcolorbox}
\section[\texorpdfstring{\kn{ಹಿಂತಿರುಗು} (hintirugu)}{ಹಿಂತಿರುಗು (hintirugu)}]{\kn{ಹಿಂತಿರುಗು} (hintirugu) — to come back, to return}

\label{lesson:KA-C278-hintirugu}



\begin{quote}
Before the new one: say the Kannada for to lend, then the Kannada for to give back.
\end{quote}

\subsection*{You'll want to know: \kn{ಹಿಂತಿರುಗು}}
\textbf{\kn{ಹಿಂತಿರುಗು}} — \emph{hintirugu} — \textquotedblleft{}to come back, to return\textquotedblright{}.

\begin{grammarlens}[title={What you've built}]
One more word for this chapter.
\end{grammarlens}

\subsection*{Your turn}
\begin{itemize}
\raggedright
  \item \emph{Say it:} \emph{hintirugu}
  \item \emph{Say it:} \emph{hintirugu}, once more
  \item \emph{Say it:} say \emph{hintirugu}
  \item \emph{Recall:} say the Kannada for to lend, then the Kannada for to give back, then say \emph{hintirugu} again
\end{itemize}

\begin{tcolorbox}[breakable,colback=teal!4,colframe=teal!35!black,title={Before you move on}]
What does \kn{ಹಿಂತಿರುಗು} mean? (\textquotedblleft{}To come back, to return\textquotedblright{}.) Say it once more.
\end{tcolorbox}
\section[\texorpdfstring{\kn{ತಡೆ} (taḍe)}{ತಡೆ (taḍe)}]{\kn{ತಡೆ} (taḍe) — to prevent, to block}

\label{lesson:KA-C278-tade}



\begin{quote}
Before the new one: say the Kannada for to give back, then the Kannada for to come back.
\end{quote}

\subsection*{You'll want to know: \kn{ತಡೆ}}
\textbf{\kn{ತಡೆ}} — \emph{taḍe} — \textquotedblleft{}to prevent, to block\textquotedblright{}.

\begin{grammarlens}[title={What you've built}]
One more word for this chapter.
\end{grammarlens}

\subsection*{Your turn}
\begin{itemize}
\raggedright
  \item \emph{Say it:} \emph{taḍe}
  \item \emph{Say it:} \emph{taḍe}, once more
  \item \emph{Say it:} say \emph{taḍe}
  \item \emph{Recall:} say the Kannada for to give back, then the Kannada for to come back, then say \emph{taḍe} again
\end{itemize}

\begin{tcolorbox}[breakable,colback=teal!4,colframe=teal!35!black,title={Before you move on}]
What does \kn{ತಡೆ} mean? (\textquotedblleft{}To prevent, to block\textquotedblright{}.) Say it once more.
\end{tcolorbox}
\section[\texorpdfstring{\kn{ತಯಾರಿಸು} (tayārisu)}{ತಯಾರಿಸು (tayārisu)}]{\kn{ತಯಾರಿಸು} (tayārisu) — to prepare, to make}

\label{lesson:KA-C278-tayarisu}



\begin{quote}
Before the new one: say the Kannada for to come back, then the Kannada for to prevent.
\end{quote}

\subsection*{You'll want to know: \kn{ತಯಾರಿಸು}}
\textbf{\kn{ತಯಾರಿಸು}} — \emph{tayārisu} — \textquotedblleft{}to prepare, to make\textquotedblright{}.

\begin{grammarlens}[title={What you've built}]
One more word for this chapter.
\end{grammarlens}

\subsection*{Your turn}
\begin{itemize}
\raggedright
  \item \emph{Say it:} \emph{tayārisu}
  \item \emph{Say it:} \emph{tayārisu}, once more
  \item \emph{Say it:} say \emph{tayārisu}
  \item \emph{Recall:} say the Kannada for to come back, then the Kannada for to prevent, then say \emph{tayārisu} again
\end{itemize}

\begin{tcolorbox}[breakable,colback=teal!4,colframe=teal!35!black,title={Before you move on}]
What does \kn{ತಯಾರಿಸು} mean? (\textquotedblleft{}To prepare, to make\textquotedblright{}.) Say it once more.
\end{tcolorbox}
